\section{\label{III_1_B}Visualisation du FRBR en graphe pour préparer la mise en place d’une nouvelle entité : la séquence et de nouvelles données produites automatiquement : alourdit les bdd donc nécessite des visualisations}

Visualisation des données et contrôle qualité (rappeller la sous partie des problèmes de navigation au sein du FRBR)

Comme nous l'avons évoqué dans la première partie, bien que le modèle FRBR soit élaboré afin d'optimiser la découvrabilité des collections, il complexifie toutefois la mise en place d'une interface intuitive et navigable, enjeu au coeur des préoccupations d'un éditeur tel que Skinsoft. L'implémentation de nouvelles fonctionnalités automatisées accentue cette complexité, avec la mise en place d'une nouvelle entité, celle du Fragment d'oeuvre, et la génération massive de données supplémentaires, qui risque de rendre moins claires les relations entre les notices. Pour ne pas entraver la lisibilité des liens entre les niveaux de description, il est nécessaire de proposer une vue d'ensemble de la structure des données au sein de S-Movies. 
Au cours de notre stage, nous avons proposé une visualisation de la structure de données sous forme de graphe relationnel. Cette proposition s'inspire notamment des travaux de la Bibliothèque Nationale Allemande qui ont amener à développer une visualisation en graphe du modèle FRBR afin de structurer les résultats de recherche indexés \footcite{zotero-item-749}. En s'appuyant sur la plateforme \textit{Culture Graph} et des outils de spatialisation tels que \textit{Gephi}, des algorithmes regroupent les œuvres sous des clés uniques et cartographient leurs liens sémantiques \footcite{zotero-item-749}. De plus, des mots-clés reliés aux entités sont également modélisés sous forme d'arrêtes \footcite{zotero-item-749}. Appliqués à une structure de données audiovisuelles, le même type de visualisation en graphe peut être imaginée, traduisant visuellement la densité des relations entre les entités du modèle. À terme, après l'implémentation des fonctionnalités automatiques d'indexation, les mots-clés générés et les times codes pourront, de la même manière, alimenter le graphe afin d'améliorer le contrôle qualité sur les données. La visualisation en graphes représente ainsi un "complément précieux" aux résultats de recherche traditionnels, permettant à l'utilisateur de comprendre beaucoup plus facilement le regroupement d'entités \footcite{zotero-item-749}. La visualisation en graphe n'est pas seulement utile pour l'amélioration de la visualisation des liens entre les entités dans les résultats de recherche, mais elle permet également d'évaluer la cohérence du modèle de données dans son ensemble. En effet, "l’impression visuelle donnée par le graphe d’une œuvre est très utile car elle attire l’ attention sur les sommets qui sont cruciaux pour la structure d’un cluster d’œuvres."\footcite{zotero-item-749}. identifier d'un coup d'œil les erreurs d'association ou les ruptures de cohérence générées par les traitements automatiques avant leur validation définitive. 
Concrètement, au sein de S-Movies, nous proposons l'intégration d'une vue d'ensemble du modèle en graphe, jouxtant la vue des résultats de recherche classique. Cette vue pourra être étendue au sien des ntoices individuelles, permettant d'obtenir une granularité différente, à partir d'un item individuel. Cette vue en graphe permettrait de faciliter la navigation entre les entités descriptives et de vérifier la cohérence des liens, dans le cadre de l'implémentation d'une indexation automatisée. 
Une description détaillée de cette proposition est présentée en annexe.

Intéropérabilité des données ?